\documentclass[11pt]{article}
\usepackage[margin=1in]{geometry}
\usepackage{amsmath, amssymb}
\usepackage{siunitx}
\setlength{\parskip}{4pt}
\setlength{\parindent}{0pt}

\title{\vspace{-2em}Higher Inertia Motors Exist in the NEM:\\
\large a Bound and What Measurements Imply}
\author{}
\date{}

\begin{document}
\maketitle
\vspace{-3em}

\paragraph{The quantity.}
The aggregate inertia constant of a composite-load-model (CLM) bus is linear in
the model's own parameters,
\begin{equation}
  H_{\mathrm{load}} \;=\; \sum_{i\in\{A,B,C\}} \frac{F_{m_i}}{\mathrm{LF}_i}\,H_i,
  \label{eq:hload}
\end{equation}
with $F_{m_i}$ the load fraction on motor class $i$, $\mathrm{LF}_i$ its loading
factor, $H_i$ its inertia constant. Only motors A, B, C carry inertia; static
(ZIP), electronic and Motor~D components contribute none.

\paragraph{The bound, inverted.}
Let $F_m=\sum_i F_{m_i}\le 1$ and $H_{\max}=\max_i H_i$. With a common loading
factor ($\mathrm{LF}=0.75$ in AEMO's set), \eqref{eq:hload} is bounded by its
largest term, $H_{\mathrm{load}}\le (F_m/\mathrm{LF})\,H_{\max}$, so a measured
load inertia demands an individual motor inertia of at least
\begin{equation}
  H_{\max} \;\ge\; \frac{\mathrm{LF}}{F_m}\;H_{\mathrm{load}}^{\mathrm{meas}}.
  \label{eq:inverse}
\end{equation}
The measured NEM/GB demand-side inertia,
$H_{\mathrm{load}}^{\mathrm{meas}}\approx\SIrange{1.4}{1.75}{\second}$, dwarfs the
$\approx\SI{0.17}{\second}$ ($\beta\approx\SI{4}{\percent}$) that \eqref{eq:hload}
returns on AEMO's set. The shortfall is in $H$ --- even at the limit $F_m\!\to\!1$, \eqref{eq:inverse} needs
$H_{\max}\gtrsim\SI{1.1}{\second}$, already above every value AEMO assigns or
recommends ($H_B=\SI{0.5}{\second}$, recommended ${\le}\,\SI{0.8}{\second}$).

\paragraph{What the gap points to.}
The composite load model parameters do not reflect load inertia observed in the NEM.
The gap points to the existence of higher inertia motors in the NEM than those represented in AEMO's set.
\end{document}
